\chapter{Conclusion}
% ============================================================

This project addresses intrusion detection for the TSN Ethernet backbone of
automotive zonal architectures, targeting attacks against TSN protocol features that
degrade timing guarantees without altering packet content or traffic volume. The
contribution is the use of protocol-native, schedule-derived signals as the detection
feature space, trained federatively across zone controllers.

At the mid-report stage the data generation path is complete and reproducible: a
four-zone TSN topology with gPTP synchronisation, three instrumented attack
configurations, fifteen seeded runs and a labelled dataset of 4{,}740 feature windows.
Delay-attack injection has been validated against its expected theoretical signature.
The five-zone ring reference topology has since been built and used to simulate three
further attack families: constant drift, asymmetric delay, and stream-redundancy attacks
on an IEEE 802.1CB protected stream (Section~\ref{sec:ring-results}). Each leaves a
measurable signature in a protocol-mechanism metric (the gPTP rate ratio, the spatial
pattern of offset shift, and the FRER elimination count) where the obvious metric is
blind or ambiguous. These are single-run observations and are not yet extracted into
the dataset.

The project's central premise has been tested ahead of schedule and holds. The test used
a proof-of-concept centralized detector, built at reduced scale for the sole purpose of
establishing feasibility. Features computed within a single zone detect the attack no
better than a majority-class predictor; two cross-zone residual features reach
0.95--0.98 macro-F1, detecting two synchronisation intervals after onset with no missed
attack windows in the best-performing arm. The gain is $+0.43$ to $+0.49$ macro-F1 and
is insensitive to model choice, which is the evidence the contribution statement
requires.

Two negative results bound the claim. Oscillator drift manipulation is not realisable as
a runtime attack under the simulation framework, which reduced the committed attack set
to two classes; a constant drift applied from the start of a run is hidden in the clock
offset but visible in the servo's rate ratio on the ring, which makes it a candidate
feature rather than a result. More significantly, no fixed threshold on the offset
residual separates attack from benign behaviour. The synchronisation convergence
transient at run start produces residuals 2.7 times larger than the attack, and because
those windows constitute roughly 11\% of benign data, no percentile-based cut-off
succeeds either. The benign class is a mixture of a convergence population and a
steady-state population with the attack lying between them in magnitude, which is what
makes a learned detector necessary rather than merely convenient, and what makes the
second residual feature, weak on its own, the one that separates the two benign
populations.

The remaining work is the detection model itself, which the proof of concept
anticipated only at reduced scale, together with the federated training pipeline, the
comparative evaluation, and extraction of the ring-topology signals, including the FRER
feature family, into the dataset. The plan sequenced a central decision gate, the
comparison between protocol-native and generic traffic features, ahead of federated
infrastructure, so that the project's primary claim was tested before further work was
built upon it; that gate has been passed.

% ============================================================
\appendix
% ============================================================